\documentclass[10pt,a4paper]{amsart}
\usepackage[margin=19mm]{geometry}
\usepackage{amsmath,amssymb,mathtools}
\usepackage[scr=rsfs,cal=euler]{mathalfa}
\usepackage{xfrac}
\usepackage[unicode,hidelinks,bookmarks=false]{hyperref}
\newcommand{\ie}{i.e.\ }
\newcommand{\cf}{cf.\ }
\newcommand{\resp}{respectively\ }
\newcommand{\Subsection}{\subsection}
\providecommand{\nobreakdash}{\nobreak\hbox{-}}
\setlength{\emergencystretch}{2em}
\multlinegap0pt
\usepackage{lmodern}
\hypersetup{
  pdftitle={A counterexample to the Foregger--Sinkhorn tie-point conjecture},
  pdfauthor={Yair Lavi}
}
\newtheorem{theorem}{Theorem}
\newtheorem{lemma}[theorem]{Lemma}
\newtheorem*{conjecture}{Tie-point conjecture}
\setlength{\emergencystretch}{3em}
\begin{document}
\title[A counterexample to the Foregger-Sinkhorn tie-point conjectu]{A counterexample to the Foregger-Sinkhorn tie-point conjecture}
\author{Yair Lavi}
\date{}
\maketitle
\noindent\emph{Source and licence.} A counterexample to the Foregger-Sinkhorn tie-point conjecture. Version arXiv:2608.13025v1 (CC BY 4.0). This material is distributed under the Creative Commons Attribution 4.0 International licence, http://creativecommons.org/licenses/by/4.0/, as recorded in the arXiv OAI licence metadata for this exact version and shown on the abstract page https://arxiv.org/abs/2608.13025v1. Full rights notice: licenses/arxiv-2608.13025-CC-BY-4.0.txt.\par\smallskip
\noindent\textit{Editorial scope.} Counted unit: the original \texttt{lemma} environment \texttt{-} at source lines 454--457 together with its complete proof at lines 459--564; the delivered document keeps source lines 449--565 so that every referenced label and every citation resolves inside it. Native editable source is the paper's own concatenated TeX; the AMS wrapper is editorial formatting only. No publisher class, upstream script or unrelated artwork is redistributed.
\section{Unique global minimum on the closed face}

We first bound the seven corners other than \(C_{111}\), and then treat
\(C_{111}\) separately.


\begin{lemma}
For every \(t\in[0,1]\), \(C_{ijk}(t)>1/50\) whenever
\((i,j,k)\ne(1,1,1)\).
\end{lemma}

\begin{proof}
Since \(x\mapsto x^5\) is strictly convex on \([0,1]\), Jensen's inequality
gives

\[
\frac{t^5+(1-t)^5}{2}
\ge\left(\frac{t+(1-t)}2\right)^5
=\frac1{32}.
\]

Hence

\[
C_{000}(t)=t^5+(1-t)^5\ge\frac1{16}>\frac1{50},         \tag{4.1}
\]

with equality in the first inequality exactly at \(t=1/2\).

For the corner \(C_{100}=C_{010}\), write

\[
C_{100}(t)=
\underbrace{\frac{t^5(2-t)}4}_{X(t)}
+\underbrace{\frac{(1-t)^4(2-t)^2}{4}}_{Y(t)}.          \tag{4.2}
\]

On \([0,1]\), \(X\) is increasing and \(Y\) is decreasing; indeed,

\[
X'(t)=\frac{t^4(5-3t)}2\ge0,
\qquad
Y'(t)=-\frac{(1-t)^3(2-t)(5-3t)}2\le0.
\]

Both \(X\) and \(Y\) are nonnegative on \([0,1]\). Three elementary estimates
now cover the interval. If \(t\le1/2\), then

\[
Y(t)\ge Y(1/2)=\frac9{256}>\frac1{50}.
\]

If \(t\ge3/5\), then

\[
X(t)\ge X(3/5)=\frac{1701}{62500}>\frac1{50}.
\]

Finally, if \(1/2\le t\le3/5\), then

\[
X(t)\ge\frac3{256}>\frac1{100},
\qquad
Y(t)\ge\frac{196}{15625}>\frac1{100},
\]

and hence \(C_{100}(t)>1/50\). The identity
\(C_{001}(t)=C_{100}(1-t)\) proves the same bound for the remaining asymmetric
corner.

For the two corners \(C_{101}=C_{011}\), exact differentiation gives

\[
\begin{aligned}
C_{101}(t)
&=\frac{(2-t)(1+t)(t^2-t+1)(7t^2-7t+2)}{16},\\
C_{101}'(t)
&=-\frac{(2t-1)(3t^2-3t+2)(7t^2-7t-8)}{16}.
\end{aligned}                                             \tag{4.3}
\]

The quadratic \(3t^2-3t+2\) is positive, while
\(7t^2-7t-8<0\), throughout \([0,1]\). Thus the derivative has the sign of
\(2t-1\), and

\[
C_{101}(t)\ge C_{101}(1/2)=\frac{27}{1024}>\frac1{50}.
\]

The last of these seven corners is

\[
\begin{aligned}
C_{110}(t)
&=\frac{(t-2)^2(7t^4-19t^3+25t^2-16t+4)}{16},\\
C_{110}'(t)
&=\frac{(t-2)(3t-2)(7t-10)(2t^2-3t+2)}{16}.
\end{aligned}                                             \tag{4.4}
\]

On \([0,1]\), the factors \(t-2\) and \(7t-10\) are negative, and
\(2t^2-3t+2\) is positive. The derivative therefore has the sign of
\(3t-2\). Consequently,

\[
C_{110}(t)\ge C_{110}(2/3)=\frac{16}{729}>\frac1{50}.
\]

We have proved

\[
C_{ijk}(t)>\frac1{50}
\qquad\text{for every }(i,j,k)\ne(1,1,1).
\]

This proves the lemma.
\end{proof}


\end{document}
